\section{Canonical products, moments, and the diagonal model}
\subsection{The index set records multiplicity}
The product index is a dependent pair. Its first component chooses an actual zero. Its second chooses one of the copies counted by that zero's analytic multiplicity.
\excerpt{FullCanonicalProduct}{multiplicity}
\excerpt{FullCanonicalProduct}{canonical_product}
For a zero of order $m$, the dependent finite type contributes exactly $m$ copies of its factor. The construction therefore does not require a global simplicity assumption. Its factor uses the reduced variable $w$.

\subsection{Absolute summability precedes the product theorem}
The weighted reciprocal-zero sum is converted to a sum over the multiplicity index.
\excerpt{FullCanonicalProduct}{product_sum}
The next source results prove that the perturbations $w/\rho$ are absolutely summable, the product is multipliable, convergence is uniform on compact sets and locally uniform, and the product defines an entire function. These are needed analytic statements. Writing an infinite product symbol does not establish them.

The product's zeros and multiplicities are then identified. Dividing the denominator by the product gives an entire, nonvanishing quotient. Growth estimates and a sublinear rigidity argument force the normalized quotient to be one. Only after those steps is the final identity short:
\excerpt{HadamardIdentification}{hadamard}
The proof depends on the product-quotient theorem and the quotient-equals-one theorem. Those in turn depend on convergence, exact multiplicities, growth reduction, and Jensen--Poisson estimates. Rows 96--102 expose these stages in the map. The earlier tail product in row 74 is one ingredient, not the complete identification.

\subsection{Recovering coefficients from zeros}
After the full product is identified, its logarithmic derivative gives a convergent Mittag--Leffler expansion away from the zeros. Rows 104--109 match derivatives and analytic Taylor coefficients with the actual multiplicity-weighted zero moments. The coefficients previously called \code{kappaCoeff} are thereby identified with zero sums.

This direction matters. The coefficient recurrence was proved algebraically before any product formula. The analytic argument does not define the denominator's coefficients by a conjectured zero expansion. It proves that the coefficients already defined equal the zero data. Tangent coefficients, reciprocal Bell formulas, and theta Bell formulas then inherit exact normalizations through named bridges.

Normalized nonzero residue sections need their own growth and zero arguments. Rows 121--132 cover their entire realization, zero-growth estimates, summability, product identification, logarithms, analytic powers, and Taylor-differential conversions. They are not inferred merely by relabeling the distinguished denominator.

\subsection{A concrete trace-class definition}
The source uses a rank-one summability definition for Hilbert trace class.
\excerpt{DiagonalNuclear}{traceclass}
The witness is a family of rank-one operators with summable products of vector norms and a sum equal to the continuous linear map. For the diagonal model, reciprocal reduced zeros form the eigenvalue family. Their absolute summability is proved from the zero theorem, then used to build the operator and its trace-class witness.

The source also proves the basis action, power traces, finite compression determinants, and their limit. The exact determinant identity is:
\excerpt{DiagonalFredholm}{fredholm}
Here \code{diagonalFredholmDeterminant} is the defined diagonal specialization, backed by convergence of finite compressions. The identity follows by matching its factors to the multiplicity-indexed canonical product. It does not assert a general Fredholm determinant theory for arbitrary operators or conceal such a theorem behind the name.

\subsection{Complete symmetric functions}
The reciprocal of the diagonal determinant has complete homogeneous symmetric-function coefficients. Rows 119--120 handle the finite combinatorics and the infinite summability argument. The latter cannot be replaced by an unqualified passage from finite-dimensional symmetric polynomials to infinitely many variables. The summable eigenvalue family is the analytic input that makes the coefficient expression meaningful.

This yields three linked descriptions of the same scalar family: formal reciprocal coefficients, zero moments and symmetric functions, and a diagonal determinant model. Each identification is a theorem with a specified source and normalization.
